% Source: https://artofproblemsolving.com/community/c3h1622939_emic_2015
% Stable identity: aops-topic-1622939-q02
Let the first number be $\overline{aa}=11a$, the second $\overline{ba}$ with $b\ne a$, $b\ge1$, and the third $b$. The number $11a$ is prime only for $a=1$.

\textbf{Case $a=1$.} Then $11$ is prime, and exactly one of $\overline{b1}$ and $b$ must be prime, for $b=2,\dots,9$:
\[
\begin{array}{c|cccccccc}
b&2&3&4&5&6&7&8&9\\\hline
\overline{b1}&21&31&41&51&61&71&81&91\\
\text{primes among }\overline{b1},b&b&\text{both}&\overline{b1}&b&\overline{b1}&\text{both}&\text{none}&\text{none}
\end{array}
\]
($51=3\cdot17$, $91=7\cdot13$.) This gives $b=2,4,5,6$: $4$ ways.

\textbf{Case $a\ge2$.} Then $11a$ is composite, so both $b$ and $\overline{ba}$ must be prime, with $b\in\{2,3,5,7\}$, $a\ge2$ and $a\ne b$:
\[
b=2:\ 23,29;\qquad b=3:\ 37;\qquad b=5:\ 53,59;\qquad b=7:\ 73,79,
\]
which is $7$ ways.

In total there are $4+7=\boxed{11}$ ways.
